% !PRINTED 317
\weberpage{317}
%<<CONT
many ways, and in particular have rendered good service in the
Tschirnhausen transformation.

We defined these functions in \S.\,4 by
\[
\begin{aligned}
f_0(x) &= 1\\
f_1(x) &= x+a_1\\
f_2(x) &= x^2+a_1x+a_2\\
&\ \cdots\cdots\cdots\cdots\cdots\cdots\cdots\cdots\\
f_{n-1}(x) &= x^{n-1}+a_1x^{n-2}+a_2x^{n-3}+\cdots+a_{n-1}\\
f(x)=f_n(x) &= x^n+a_1x^{n-1}+a_2x^{n-2}+\cdots+a_{n-1}x+a_n.
\end{aligned}
\]
and for them set up the recursion formula
\[
f_v(x)-xf_{v-1}(x)=a_v.
\]
One now has, as has already been shown in the place cited,
\weq{3}{f(x)=(x-\alpha)\left[x^{n-1}f_0(\alpha)+x^{n-2}f_1(\alpha)+\cdots+f_{n-1}(\alpha)\right]+f(\alpha).}
A glance at this formula shows that if $\alpha$ is a number which makes
the functions
\weq{4}{f_0(\alpha),\ f_1(\alpha),\ f_2(\alpha)\ldots f_{n-1}(\alpha),\ f(\alpha)}
positive, then no positive $x$ greater than $\alpha$ can exist which
makes $f(x)$ zero, since in that case nothing but positive terms stand
on the right-hand side of (3).

If, then, a positive $\alpha$ is so determined that the functions
(4) are all positive, this is an upper limit for the positive
roots of $f(x)$.

In the same way one obtains a lower limit $-\beta$ for the negative roots,
if one determines the positive quantity $\beta$ so that
\weq{5}{f_0(-\beta),\ -f_1(-\beta),\ f_2(-\beta),\ \ldots (-1)^n f(-\beta)}
all become positive.

This determination of the limits is indeed simpler to handle in
calculation, but under some circumstances it gives less accurate limits than
the \gesperrt{Newton} determination.

For if one differentiates (3) repeatedly with respect to $x$ and then
sets $x=\alpha$, one obtains, under the assumption that the
quantities (4) are all positive, for the derivatives $f'(x)f''(x)\ldots$
expressions which for $x=\alpha$ consist of nothing but positive terms.
If, then, the expressions (4) are positive, the derivatives (1) are
also necessarily all positive, while the converse is not
necessarily the case.
